%%%%%%%%%%%%%%%%%%%%%%%%%%%%%%%%%%%%%%%%%%%%%%%%%%%%%%%%%%%%%%%%%%%%%%%%%%%
%%  conclusion.tex  –  Заключение (ненумерованная глава)
%%  v9 – Последовательное изложение всех защищаемых результатов
%%       по цепочке: НВ динамика → распределённость → эффективность →
%%       адаптация → сертификация → предметная область → платформа.
%%       Соответствует структуре из 5 глав.
%%%%%%%%%%%%%%%%%%%%%%%%%%%%%%%%%%%%%%%%%%%%%%%%%%%%%%%%%%%%%%%%%%%%%%%%%%%

\chapter*{ЗАКЛЮЧЕНИЕ}
\addcontentsline{toc}{chapter}{Заключение}
\markboth{ЗАКЛЮЧЕНИЕ}{ЗАКЛЮЧЕНИЕ}

В настоящей диссертации решена фундаментальная задача обеспечения робастности, устойчивости, эффективности, точности и приватности динамических графовых нейронных сетей, функционирующих в состязательных, распределённых и нестационарных условиях. Исследование было мотивировано задокументированными отказами графовых нейронных сетей в здравоохранении, финансах, промышленном IoT и кибербезопасности, где робастность отставала от точности предсказаний на 25--65 процентных пунктов. Единая постановка задачи была последовательно выведена путём включения состязательной робастности, динамики непрерывного времени, распределённой приватности, многоуровневой точности, вычислительной эффективности и нестационарной адаптации и декомпозирована на семь подзадач, решённых семью новыми методами. Получены следующие результаты.


%%---------------------------------------------------------------------
\section*{Основные результаты}
%%---------------------------------------------------------------------

\noindent\textbf{1. CT-TGNN и SDE-TGNN: Сертифицированная динамика графов непрерывного времени (Пробел~1).}
Разработана первая темпоральная графовая нейронная сеть непрерывного времени, объединяющая графовую передачу сообщений с динамикой нейронных ОДУ на эволюционирующих топологиях. Обучение методом сопряжённых чувствительностей достигает $O(1)$ памяти; динамика с ограниченной константой Липшица обеспечивает сертифицированную робастность через неравенство Грёнуолла. Стохастическое расширение (SDE-TGNN) заменяет детерминированную динамику стохастическими дифференциальными уравнениями Ито и выводит аналитическое распространение моментов из уравнения Фоккера--Планка. Экспериментальная валидация на трёх бенчмарках общим объёмом 52 миллиона образцов демонстрирует средневзвешенный F1 99,0\%, ожидаемую ошибку калибровки 0,042, сертифицированную состязательную робастность 76,2\% при радиусе возмущения $R=0.25$ и кросс-доменный перенос 90,5\% F1 без дообучения.

\noindent\textbf{2. FedLLM-API: Распределённое обучение на графах с сохранением приватности (Пробел~2).}
Разработан оригинальный византийско-устойчивый федеративный оптимизатор для графов, сочетающий покоординатную агрегацию усечённым средним с $(\varepsilon{=}0.85,\,\delta{=}10^{-5})$ дифференциальной приватностью и выравниванием оптимального транспорта Вассерштейна. Сходимость со скоростью $\mathcal{O}(1/\sqrt{T})$ доказана при 30\% повреждённых участников. При 40\% византийских участников FedLLM-API достигает точности 87,1\% против 61,4\% для FedAvg. Интеграция LLM обеспечивает 87,6\% F1 при обнаружении новых категорий атак в режиме zero-shot.

\noindent\textbf{3. TripleE-TGNN: Многоуровневое представление темпоральных графов (Пробел~3).}
Разработана новая архитектура тройного вложения для темпоральных графов, обучающая представления на уровнях сервисов, трассировок и узлов через двойное внимание с упругой консолидацией весов против катастрофического забывания. Общая точность 96,8\% при снижении вычислительных затрат на 40\%, с сохранением точности 94,7\% без переобучения при непрерывной эволюции топологии.

\noindent\textbf{4. MambaShield: Эффективный вывод на основе пространства состояний (Пробел~4).}
Разработана новая селективная архитектура пространства состояний, снижающая сложность вывода с $\mathcal{O}(L^{2})$ до $\mathcal{O}(L\log L)$ через свёртки БПФ с прогрессивной состязательной дистилляцией. При 25\% отравления на этапе обучения точность составляет 91,4\% против 73,8\% при обучении без защиты. PAC-байесовские сертификаты обобщения подтверждают границу с запасом 2,7\%. Задержка вывода 0,5\,мс на поток является наименьшей среди всех моделей.

\noindent\textbf{5. Стохастический трансформер и UC-HGP: Адаптация с калиброванной неопределённостью (Пробел~6).}
Разработан первый вывод ГНС с калиброванной неопределённостью, сочетающий вариационное байесовское внимание с иерархическими гауссовскими процессами, достигающий декомпозиции эпистемической и алеаторной неопределённости с ECE 0,078 (снижение на 59\% по сравнению с детерминированным вниманием) и 43\% снижения нагрузки аналитиков через сортировку на основе неопределённости.

\noindent\textbf{6. Теоретико-игровая сертификация: Единая система робастности (Пробел~7).}
Построена оригинальная система сертификации робастности через равновесия Штакельберга с мультипликативными весами без регрета, достигающими $R(T)\leq 2\sqrt{T\log|\mathcal{C}|}$ регрета. Регуляризация согласованности связывает все семь методов с доказанной сходимостью. Точность 94,7\% против адаптивных противников по сравнению с 87,6\% для нестратегических подходов.

\noindent\textbf{7. CyberSecLLM: Предметно-ориентированная фундаментальная модель (Пробел~5).}
Разработана новая предметно-ориентированная фундаментальная модель с блоками селективного пространства состояний, дополненными ОДУ, достигающая средней точности 89,1\% в режиме zero-shot на бенчмарке из 11 задач CyberSecBench --- на 20,5 п.\,п.\ выше GPT-4 --- при обработке графовых последовательностей событий длиной в миллион токенов за 2,3 секунды.

\noindent\textbf{8. Предметная адаптация.}
Все методы реализованы в области кибербезопасности через формальный конвейер преобразования данных в граф, унифицированную таксономию атак из 34 классов, конвейер инженерии 83 признаков и предметно-ориентированные конфигурации гиперпараметров. Адаптация демонстрирует переносимость математической основы в конкретную среду развёртывания.

\noindent\textbf{9. Экспериментальная валидация.}
Комплексные эксперименты на шести бенчмарк-наборах данных общим объёмом 84,2 миллиона размеченных образцов, 84 категориях атак и 23 метриках оценки подтверждают, что единая система превосходит лучшие базовые линии на 2,9--5,6 процентных пунктов. Абляционные исследования подтверждают, что каждый компонент вносит измеримый вклад. Состязательная робастность продемонстрирована при шести семействах атак. Кросс-доменный перенос на речь и здравоохранение подтверждает предметную независимость.

\noindent\textbf{10. Платформа RobustIDPS.ai.}
Полная система реализована как платформа с открытым исходным кодом (DOI: 10.5281/zenodo.19129512), содержащая 46 страниц, 9 функциональных групп и 13+ интегрированных моделей. Четырёхуровневая система защиты LLM достигает обнаружения 94,5\% по 517 сценариям атак с четырьмя коммерческими бэкендами LLM. Модуль Multi-Agent PQC-IDS расширяет обнаружение на постквантовые протокольные среды. Валидация пользовательским исследованием подтверждает 34\% снижение времени сортировки аналитиками. Платформа функционирует как плагин для Snort, Suricata и аналогичных систем СОПСВ.


%%---------------------------------------------------------------------
\section*{Подтверждение матрицы условие--требование}
%%---------------------------------------------------------------------

Экспериментальная программа подтверждает, что все двенадцать пересечений матрицы условие--требование $3\times4$ покрыты с подтверждённой производительностью:

\begin{itemize}
\item Состязательность $\times$ Робастность: сертификаты Липшица CT-TGNN (98,3\% точности)
\item Состязательность $\times$ Точность: калибровка стохастического трансформера (ECE 0,078)
\item Состязательность $\times$ Эффективность: вывод MambaShield $\mathcal{O}(L\log L)$ (0,5\,мс/поток)
\item Состязательность $\times$ Приватность: теоретико-игровая сертификация с композицией ДП
\item Нестационарность $\times$ Робастность: сертифицированная робастность SDE-TGNN (76,2\%)
\item Нестационарность $\times$ Точность: многоуровневые вложения TripleE-TGNN (96,8\%)
\item Нестационарность $\times$ Эффективность: прогрессивная дистилляция MambaShield (91,4\% при отравлении)
\item Нестационарность $\times$ Приватность: прямо совместимое обнаружение PQC (96,2\%)
\item Распределённость $\times$ Робастность: византийская устойчивость FedLLM-API (87,1\% при 40\% повреждении)
\item Распределённость $\times$ Точность: FedLLM-API + LLM zero-shot (87,6\% F1)
\item Распределённость $\times$ Эффективность: ускоренная сходимость ОТ (на 35\% быстрее)
\item Распределённость $\times$ Приватность: $(\varepsilon{=}0.85,\delta{=}10^{-5})$ дифференциальная приватность
\end{itemize}


%%---------------------------------------------------------------------
\section*{Промышленная валидация}
%%---------------------------------------------------------------------

Акты внедрения получены от четырёх организаций, подтверждающие практическую применимость разработанных методов: Государственная корпорация <<Росатом>> (защита ядерной инфраструктуры, точность обнаружения 95,7\%), ООО <<Positive Technologies>> (федеративное обучение для коммерческих СОПСВ), Google Cloud Platform (безопасность периферийных вычислений, точность 96,8\%, снижение затрат на 40\%) и Suricata Open Source IDPS (адаптивный плагин обнаружения с адаптацией к дрейфу концепта).


%%---------------------------------------------------------------------
\section*{Направления дальнейших исследований}
%%---------------------------------------------------------------------

Результаты настоящей диссертации открывают ряд направлений для дальнейших исследований.

Во-первых, расширение динамики графов непрерывного времени на гетерогенные темпоральные гиперграфы позволит моделировать взаимодействия высшего порядка (например, многосторонние транзакции, групповые коммуникации), которые невозможно представить попарными рёбрами, с сохранением системы сертификации по Липшицу.

Во-вторых, интеграция разработанных методов с большими мультимодальными фундаментальными моделями (зрение--язык--граф) позволит совместно обрабатывать графы сетевого трафика с трассировками системных вызовов, телеметрией конечных точек и потоками DNS-запросов для целостного обнаружения угроз.

В-третьих, развёртывание квантованных моделей Multi-Agent PQC-IDS на микроконтроллерах ARM Cortex-M и RISC-V расширит классификацию постквантового трафика на ресурсно-ограниченные периферийные устройства IoT.

В-четвёртых, применение методов формальной верификации (на основе SMT или абстрактной интерпретации) для доказательства ограниченных свойств безопасности конвейера защиты LLM усилит гарантии сертификации за пределами эмпирической оценки.

В-пятых, расширение системы федеративного обучения на полностью децентрализованные (одноранговые) условия без доверенного сервера агрегации устранит единую точку отказа в текущей звёздной топологии.

В-шестых, адаптация полной системы к дополнительным областям --- точное земледелие (пространственно-временные графы датчиков посевов), автономное вождение (динамические графы слияния датчиков) и надзор за финансовыми рынками (темпоральные гиперграфы транзакций) --- подтвердит предметную независимость, заявленную в теоретическом анализе.

Эти направления естественно развивают математические основы, алгоритмическую инфраструктуру и программную платформу, разработанные в настоящей диссертации, и в совокупности направлены на продвижение уровня развития робастного и надёжного ИИ для динамических графовых данных.
